% =====================================================================
%  libreta.tex
%  ========================
%
%
%  Descripción:
%  ------------
%  Apuntes de clase del curso de Matemáticas para Computación.
%
%
%  Metadata:
%  ----------
%   * Author : zxxz6 (Bryan Violante Arriaga)
%   * Version: 2.1.0
%
%
%  History:
%  ------------
%  Author      Date            Description
%  zxxz6       26/08/2026      Pase a LaTeX la toma cruda del 26/08
%  zxxz6       24/08/2026      Pase a LaTeX la toma cruda del 24/08
%  zxxz6       19/08/2026      Pase a LaTeX la toma cruda del 19/08
%  zxxz6       19/08/2026      Creation
%
%
% =====================================================================
\documentclass[11pt,letterpaper]{article}

% =====================================================================
%  DATOS DE LA MATERIA
% =====================================================================
\newcommand{\materia}{Matemáticas para Computación}
\newcommand{\profesor}{Dra. Hayde Peregrina Barreto \\
                       Dr. Luis Villaseñor Pineda}
\newcommand{\periodo}{Otoño 2026}
\newcommand{\alumno}{Bryan Ezequiel Violante Arriaga}

% =====================================================================
%  Preambulo.tex
%  ========================
%
%
%  Descripción:
%  ------------
%  Preámbulo compartido de los apuntes de clase: paquetes, estilo de
%  página, entornos de teorema, las cuatro cajas de color, el estilo de
%  listings y el pseudocódigo en español. Se carga con \input desde
%  cada documento en lugar de copiarlo, para que un arreglo aquí llegue
%  a todas las libretas de una vez.
%
%
%  Considerations:
%  ------------
%  - Se carga con \input y ruta relativa, no con \usepackage, porque
%    \usepackage solo busca en la carpeta del documento y en el árbol
%    de TeX. Un .sty instalado en TEXMFHOME evitaría la ruta, pero vive
%    fuera del repo y una clona nueva no compilaría
%  - Los cuatro datos de la materia se declaran en el documento ANTES
%    del \input. hyperref expande \materia al cargarse, así que si no
%    existe todavía la compilación truena
%  - Cuando la materia la imparte más de un profesor, los nombres van
%    en un solo \profesor separados por \\. La portadilla los imprime
%    dentro de un center, que es donde \\ sí corta renglón. Repetir
%    \newcommand{\profesor} no es opción: truena con "already defined"
%  - algorithmicx no tiene opción de idioma y babel no lo traduce, así
%    que las palabras clave del pseudocódigo se renombran a mano
%  - Los nombres de los entornos van sin acento (\begin{definicion})
%    aunque el rótulo impreso sí lo lleve
%  - babel se carga con es-noshorthands porque en español el " queda
%    activo y se come el texto que le sigue, además de romper listings
%
%
%  Metadata:
%  ----------
%   * Author : zxxz6 (Bryan Violante Arriaga)
%   * Version: 1.7.0
%
%
%  History:
%  ------------
%  Author      Date            Description
%  zxxz6       24/09/2026      El nombre del alumno en el encabezado
%  zxxz6       24/09/2026      Encabezado de una hoja de un ejercicio
%  zxxz6       23/09/2026      Macros de los documentos de soluciones
%  zxxz6       30/08/2026      Macros de lim, limsup y liminf sobre n
%  zxxz6       25/08/2026      Macros de Omega, Theta, o y omega chicas
%  zxxz6       20/08/2026      Cargue pgfplots para graficas de datos
%  zxxz6       20/08/2026      Cargue tikz para diagramas de conjuntos
%  zxxz6       19/08/2026      Documente \profesor con varios nombres
%  zxxz6       19/08/2026      Agregue la caja Idea Random
%  zxxz6       18/08/2026      Creation
%
%
% =====================================================================
%  USO
%  ---
%  \documentclass[11pt,letterpaper]{article}
%  \newcommand{\materia}{...}      % los cuatro datos van arriba
%  \newcommand{\profesor}{...}
%  \newcommand{\periodo}{...}
%  \newcommand{\alumno}{...}
%  % =====================================================================
%  Preambulo.tex
%  ========================
%
%
%  Descripción:
%  ------------
%  Preámbulo compartido de los apuntes de clase: paquetes, estilo de
%  página, entornos de teorema, las cuatro cajas de color, el estilo de
%  listings y el pseudocódigo en español. Se carga con \input desde
%  cada documento en lugar de copiarlo, para que un arreglo aquí llegue
%  a todas las libretas de una vez.
%
%
%  Considerations:
%  ------------
%  - Se carga con \input y ruta relativa, no con \usepackage, porque
%    \usepackage solo busca en la carpeta del documento y en el árbol
%    de TeX. Un .sty instalado en TEXMFHOME evitaría la ruta, pero vive
%    fuera del repo y una clona nueva no compilaría
%  - Los cuatro datos de la materia se declaran en el documento ANTES
%    del \input. hyperref expande \materia al cargarse, así que si no
%    existe todavía la compilación truena
%  - Cuando la materia la imparte más de un profesor, los nombres van
%    en un solo \profesor separados por \\. La portadilla los imprime
%    dentro de un center, que es donde \\ sí corta renglón. Repetir
%    \newcommand{\profesor} no es opción: truena con "already defined"
%  - algorithmicx no tiene opción de idioma y babel no lo traduce, así
%    que las palabras clave del pseudocódigo se renombran a mano
%  - Los nombres de los entornos van sin acento (\begin{definicion})
%    aunque el rótulo impreso sí lo lleve
%  - babel se carga con es-noshorthands porque en español el " queda
%    activo y se come el texto que le sigue, además de romper listings
%
%
%  Metadata:
%  ----------
%   * Author : zxxz6 (Bryan Violante Arriaga)
%   * Version: 1.7.0
%
%
%  History:
%  ------------
%  Author      Date            Description
%  zxxz6       24/09/2026      El nombre del alumno en el encabezado
%  zxxz6       24/09/2026      Encabezado de una hoja de un ejercicio
%  zxxz6       23/09/2026      Macros de los documentos de soluciones
%  zxxz6       30/08/2026      Macros de lim, limsup y liminf sobre n
%  zxxz6       25/08/2026      Macros de Omega, Theta, o y omega chicas
%  zxxz6       20/08/2026      Cargue pgfplots para graficas de datos
%  zxxz6       20/08/2026      Cargue tikz para diagramas de conjuntos
%  zxxz6       19/08/2026      Documente \profesor con varios nombres
%  zxxz6       19/08/2026      Agregue la caja Idea Random
%  zxxz6       18/08/2026      Creation
%
%
% =====================================================================
%  USO
%  ---
%  \documentclass[11pt,letterpaper]{article}
%  \newcommand{\materia}{...}      % los cuatro datos van arriba
%  \newcommand{\profesor}{...}
%  \newcommand{\periodo}{...}
%  \newcommand{\alumno}{...}
%  % =====================================================================
%  Preambulo.tex
%  ========================
%
%
%  Descripción:
%  ------------
%  Preámbulo compartido de los apuntes de clase: paquetes, estilo de
%  página, entornos de teorema, las cuatro cajas de color, el estilo de
%  listings y el pseudocódigo en español. Se carga con \input desde
%  cada documento en lugar de copiarlo, para que un arreglo aquí llegue
%  a todas las libretas de una vez.
%
%
%  Considerations:
%  ------------
%  - Se carga con \input y ruta relativa, no con \usepackage, porque
%    \usepackage solo busca en la carpeta del documento y en el árbol
%    de TeX. Un .sty instalado en TEXMFHOME evitaría la ruta, pero vive
%    fuera del repo y una clona nueva no compilaría
%  - Los cuatro datos de la materia se declaran en el documento ANTES
%    del \input. hyperref expande \materia al cargarse, así que si no
%    existe todavía la compilación truena
%  - Cuando la materia la imparte más de un profesor, los nombres van
%    en un solo \profesor separados por \\. La portadilla los imprime
%    dentro de un center, que es donde \\ sí corta renglón. Repetir
%    \newcommand{\profesor} no es opción: truena con "already defined"
%  - algorithmicx no tiene opción de idioma y babel no lo traduce, así
%    que las palabras clave del pseudocódigo se renombran a mano
%  - Los nombres de los entornos van sin acento (\begin{definicion})
%    aunque el rótulo impreso sí lo lleve
%  - babel se carga con es-noshorthands porque en español el " queda
%    activo y se come el texto que le sigue, además de romper listings
%
%
%  Metadata:
%  ----------
%   * Author : zxxz6 (Bryan Violante Arriaga)
%   * Version: 1.7.0
%
%
%  History:
%  ------------
%  Author      Date            Description
%  zxxz6       24/09/2026      El nombre del alumno en el encabezado
%  zxxz6       24/09/2026      Encabezado de una hoja de un ejercicio
%  zxxz6       23/09/2026      Macros de los documentos de soluciones
%  zxxz6       30/08/2026      Macros de lim, limsup y liminf sobre n
%  zxxz6       25/08/2026      Macros de Omega, Theta, o y omega chicas
%  zxxz6       20/08/2026      Cargue pgfplots para graficas de datos
%  zxxz6       20/08/2026      Cargue tikz para diagramas de conjuntos
%  zxxz6       19/08/2026      Documente \profesor con varios nombres
%  zxxz6       19/08/2026      Agregue la caja Idea Random
%  zxxz6       18/08/2026      Creation
%
%
% =====================================================================
%  USO
%  ---
%  \documentclass[11pt,letterpaper]{article}
%  \newcommand{\materia}{...}      % los cuatro datos van arriba
%  \newcommand{\profesor}{...}
%  \newcommand{\periodo}{...}
%  \newcommand{\alumno}{...}
%  \input{../../templates/Preambulo.tex}
%  \begin{document}
%  \portadilla
%
%  Con dos o mas profesores se declara UN solo \profesor y los nombres
%  se separan con \\, uno por renglon:
%
%  \newcommand{\profesor}{Dra. Fulana de Tal \\
%                         Dr. Mengano de Cual}
% =====================================================================


% ==================== IDIOMA Y CODIFICACIÓN ====================
\usepackage[utf8]{inputenc}
\usepackage[T1]{fontenc}
\usepackage[spanish,es-tabla,es-noshorthands]{babel}
\usepackage{lmodern}
\usepackage{microtype}

% ==================== PÁGINA ====================
\usepackage[letterpaper,margin=2.3cm,headheight=14pt]{geometry}
\usepackage{parskip}
\usepackage{fancyhdr}
\usepackage{enumitem}

% ==================== MATEMÁTICAS ====================
\usepackage{amsmath,amssymb,amsthm,mathtools}

% ==================== FIGURAS, TABLAS, CÓDIGO ====================
\usepackage{graphicx}
\usepackage{booktabs}
\usepackage{array}
% float da el especificador [H], que clava una tabla donde se
% escribio. Una traza que LaTeX manda dos paginas adelante deja
% de leerse junto al paso que la produjo
\usepackage{float}
\usepackage{xcolor}
% tikz ya viene arrastrado por tcolorbox[most], pero se carga
% explicito para no depender de ese detalle; la libreria babel
% protege a tikz de los caracteres activos del español
\usepackage{tikz}
\usetikzlibrary{babel}
% pgfplots dibuja graficas de datos (ejes, escalas log) sobre tikz
\usepackage{pgfplots}
\pgfplotsset{compat=1.18}
\usepackage{listings}
\usepackage[ruled,section]{algorithm}
\usepackage{algpseudocode}
\usepackage[most]{tcolorbox}
% soul da \hl, el unico resaltado que parte renglon; \colorbox no
\usepackage{soul}
\usepackage{hyperref}

% =====================================================================
%  ESTILO
% =====================================================================

% --- encabezado y pie de cada página ---
\pagestyle{fancy}
\fancyhf{}
\fancyhead[L]{\small\textsc{\materia}}
\fancyhead[R]{\small\periodo}
\fancyfoot[C]{\small\thepage}
\renewcommand{\headrulewidth}{0.4pt}

% --- enlaces ---
\hypersetup{
  colorlinks=true,
  linkcolor=black,
  urlcolor=blue,
  citecolor=black,
  pdftitle={Apuntes de \materia},
  pdfauthor={\alumno}
}

% --- paleta ---
\definecolor{acento}{HTML}{1F4E79}
\definecolor{fondoIdea}{HTML}{EAF2FB}
\definecolor{fondoOjo}{HTML}{FDF0E6}
\definecolor{fondoPend}{HTML}{F2F0F7}
\definecolor{comentario}{HTML}{5A7D5A}
\definecolor{cadena}{HTML}{A03030}
\definecolor{fondoCodigo}{HTML}{F7F7F7}
\definecolor{fondoResultado}{HTML}{EAF5EC}
\definecolor{marcador}{HTML}{FFF3A3}

% --- entornos de teorema; numeran por sección, o sea por clase ---
\theoremstyle{definition}
\newtheorem{definicion}{Definición}[section]
\newtheorem{ejemplo}[definicion]{Ejemplo}
\newtheorem{ejercicio}[definicion]{Ejercicio}

\theoremstyle{plain}
\newtheorem{teorema}[definicion]{Teorema}
\newtheorem{lema}[definicion]{Lema}
\newtheorem{corolario}[definicion]{Corolario}
\newtheorem{proposicion}[definicion]{Proposición}

\theoremstyle{remark}
\newtheorem*{observacion}{Observación}
\newtheorem*{quick_sum}{Resumen rápido}

% --- cajas de color para lo que se relee antes del examen ---
\newtcolorbox{idea}[1][]{
  colback=fondoIdea, colframe=acento, boxrule=0.6pt,
  arc=2pt, left=6pt, right=6pt, top=5pt, bottom=5pt,
  title={Idea clave}, fonttitle=\bfseries\small, #1
}
\newtcolorbox{ojo}[1][]{
  colback=fondoOjo, colframe=orange!70!black, boxrule=0.6pt,
  arc=2pt, left=6pt, right=6pt, top=5pt, bottom=5pt,
  title={Ojito...}, fonttitle=\bfseries\small, #1
}
\newtcolorbox{pendiente}[1][]{
  colback=fondoPend, colframe=violet!60!black, boxrule=0.6pt,
  arc=2pt, left=6pt, right=6pt, top=5pt, bottom=5pt,
  title={Pendiente por resolver}, fonttitle=\bfseries\small, #1
}
\newtcolorbox{idear}[1][]{
  colback=fondoPend, colframe=pink!60!black, boxrule=0.6pt,
  arc=2pt, left=6pt, right=6pt, top=5pt, bottom=5pt,
  title={Idea Random}, fonttitle=\bfseries\small, #1
}

% --- el resultado al que llego un ejercicio, para que se vea de
%     lejos al hojear la hoja antes del examen ---
\newtcolorbox{resultado}[1][]{
  colback=fondoResultado, colframe=green!45!black, boxrule=0.6pt,
  arc=2pt, left=6pt, right=6pt, top=5pt, bottom=5pt,
  title={Resultado}, fonttitle=\bfseries\small, #1
}

% --- código ---
\lstdefinestyle{codigo}{
  backgroundcolor=\color{fondoCodigo},
  basicstyle=\ttfamily\footnotesize,
  keywordstyle=\color{acento}\bfseries,
  commentstyle=\color{comentario}\itshape,
  stringstyle=\color{cadena},
  numbers=left, numberstyle=\tiny\color{gray}, numbersep=8pt,
  frame=single, rulecolor=\color{gray!40},
  breaklines=true, showstringspaces=false,
  tabsize=4, captionpos=b,
  extendedchars=true, inputencoding=utf8,
  literate={á}{{\'a}}1 {é}{{\'e}}1 {í}{{\'i}}1
           {ó}{{\'o}}1 {ú}{{\'u}}1 {ñ}{{\~n}}1
           {Ü}{{\"U}}1 {ü}{{\"u}}1 {¿}{{?`}}1
           {¡}{{!`}}1
}
\lstset{style=codigo}
\renewcommand{\lstlistingname}{Código}
\renewcommand{\lstlistlistingname}{Índice de código}
% --- pseudocódigo; algorithmicx arma el cuerpo y algorithm el
%     flotante que lo numera y lo pone en el índice. Las palabras
%     clave vienen en inglés de fábrica y babel no las toca, así que
%     se renombran una por una ---
\floatname{algorithm}{Algoritmo}
\renewcommand{\listalgorithmname}{Índice de algoritmos}

\algrenewcommand\algorithmicrequire{\textbf{Entrada:}}
\algrenewcommand\algorithmicensure{\textbf{Salida:}}
\algrenewcommand\algorithmicend{\textbf{fin}}
\algrenewcommand\algorithmicif{\textbf{si}}
\algrenewcommand\algorithmicthen{\textbf{entonces}}
\algrenewcommand\algorithmicelse{\textbf{si no}}
\algrenewcommand\algorithmicfor{\textbf{para}}
\algrenewcommand\algorithmicforall{\textbf{para todo}}
\algrenewcommand\algorithmicdo{\textbf{hacer}}
\algrenewcommand\algorithmicwhile{\textbf{mientras}}
\algrenewcommand\algorithmicrepeat{\textbf{repetir}}
\algrenewcommand\algorithmicuntil{\textbf{hasta que}}
\algrenewcommand\algorithmicloop{\textbf{ciclo}}
\algrenewcommand\algorithmicprocedure{\textbf{procedimiento}}
\algrenewcommand\algorithmicfunction{\textbf{función}}
\algrenewcommand\algorithmicreturn{\textbf{devolver}}
\algrenewcomment[1]{\hfill{\color{comentario}\itshape // #1}}

% Las dos que algorithmicx no trae. El \ del final es obligatorio: sin
% el TeX se come el espacio y sale "hastan-1" en vez de "hasta n-1"
\algnewcommand\Hasta{\textbf{hasta}\ }
\algnewcommand\Intercambiar{\textbf{intercambiar}\ }

% --- abreviaturas que se usan a cada rato ---
\newcommand{\N}{\mathbb{N}}
\newcommand{\Z}{\mathbb{Z}}
\newcommand{\Q}{\mathbb{Q}}
\newcommand{\R}{\mathbb{R}}
\newcommand{\C}{\mathbb{C}}
\newcommand{\abs}[1]{\left\lvert #1 \right\rvert}
\newcommand{\norm}[1]{\left\lVert #1 \right\rVert}
\newcommand{\set}[1]{\left\{ #1 \right\}}
\newcommand{\Ord}[1]{\mathcal{O}\!\left(#1\right)}
% Mayuscula la notacion grande, minuscula la chica
\newcommand{\Omg}[1]{\Omega\!\left(#1\right)}
\newcommand{\Tht}[1]{\Theta\!\left(#1\right)}
\newcommand{\ord}[1]{o\!\left(#1\right)}
\newcommand{\omg}[1]{\omega\!\left(#1\right)}
% Los tres limites sobre n, que salen en cada renglon del criterio
% del cociente para la notacion asintotica
\newcommand{\LM}{\lim_{n \to \infty}}
\newcommand{\LS}{\limsup_{n \to \infty}}
\newcommand{\LI}{\liminf_{n \to \infty}}

% --- encabezado de cada sesión: título a la izquierda, fecha a la
%     derecha. El argumento corto es el que sale en el índice ---
\newcommand{\clase}[2]{%
  \section[#1]{#1 \hfill \normalfont\small\textit{#2}}%
}

% --- portadilla y tabla de contenidos. Se llama una vez, justo
%     despues de \begin{document}. El renglon del profesor va dentro
%     del center, asi que \profesor puede traer varios nombres
%     separados por \\ y salen apilados y centrados sin tocar nada ---
\newcommand{\portadilla}{%
  \begin{center}
    {\LARGE\bfseries\materia\par}
    \vspace{0.4em}
    {\large \profesor \par}
    \vspace{0.2em}
    {\small \periodo\quad-\quad\alumno\par}
    \vspace{0.3em}
    \rule{\linewidth}{0.5pt}
  \end{center}
  \vspace{0.5em}
  \tableofcontents
  \vspace{1em}
  \rule{\linewidth}{0.5pt}%
}

% =====================================================================
%  DOCUMENTOS DE SOLUCIONES
%  Lo que usa una tanda de ejercicios resueltos. Una libreta no toca
%  nada de aqui, pero vive en el preambulo compartido para no tener
%  dos archivos que arreglar cuando cambie el estilo.
% =====================================================================

% --- encabezado compacto: titulo, subtitulo y regla. Sustituye a
%     \portadilla cuando el documento es una tanda de ejercicios: sin
%     indice, porque se hojea de corrido, no se navega ---
\newcommand{\encabezadoSoluciones}[2]{%
  % El paquete algorithm se carga con la opcion [section] para que en
  % una libreta el pseudocodigo se numere por clase. Aqui no hay
  % secciones, asi que esa numeracion sale como "Algoritmo 0.1". Se
  % corrige al vuelo, y solo en los documentos que llaman a este
  % encabezado: la libreta no se entera.
  \renewcommand{\thealgorithm}{\arabic{algorithm}}%
  \begin{center}
    {\LARGE\bfseries #1\par}
    \vspace{0.3em}
    {\small\itshape #2\par}
    \vspace{0.5em}
    \rule{\linewidth}{0.5pt}
  \end{center}
  \vspace{0.3em}%
}

% --- hoja de un solo ejercicio: un renglon chico arriba con el
%     titulo y su numero, y ya. Sin titulo grande y sin pie de pagina,
%     porque en una hoja de dos paginas cada centimetro que se va en
%     adornos es un paso del desarrollo que no cabe ---
\newcommand{\encabezadoEjercicio}[3]{%
  % Misma correccion que en \encabezadoSoluciones: sin secciones, el
  % pseudocodigo saldria numerado "0.1"
  \renewcommand{\thealgorithm}{\arabic{algorithm}}%
  \fancyhf{}%
  \fancyhead[L]{\small\textsc{#1} (ejercicio #2 de #3)}%
  % El nombre va a la derecha porque las hojas de respuestas se
  % entregan y el profesor pide que cada una lo lleve
  \fancyhead[R]{\small\alumno}%
  \renewcommand{\headrulewidth}{0.4pt}%
  \pagestyle{fancy}%
  \thispagestyle{fancy}%
}

% --- bloque tematico, numerado con romanos: I, II, III. Agrupa los
%     ejercicios que comparten tema. No dibuja regla: la del
%     \problema que sigue ya separa el titulo del primer ejercicio,
%     y dos rayas seguidas dejan un hueco feo ---
\newcounter{parte}
\renewcommand{\theparte}{\Roman{parte}}
\newcommand{\parte}[1]{%
  \bigskip
  \refstepcounter{parte}%
  \noindent{\large\bfseries \theparte. #1\par}%
}

% --- un ejercicio resuelto. Se numera solo y corrido a lo largo del
%     documento, sin reiniciar en cada parte, para poder decir "el 14"
%     sin aclarar de que bloque ---
\newcounter{problema}
\newcommand{\problema}[1]{%
  \bigskip
  \noindent\rule{\linewidth}{0.4pt}\par
  \vspace{0.3em}
  \refstepcounter{problema}%
  \noindent{\bfseries Ejercicio \theproblema\ --- #1\par}
  \vspace{0.3em}%
}

% --- rotulo del paso que sigue: Caso base, Hipotesis inductiva, Paso
%     inductivo. El texto sigue en el mismo renglon ---
\newcommand{\paso}[1]{\par\vspace{0.3em}\noindent\textbf{#1}\ }

% --- la respuesta de un ejercicio, centrada y en un cuadro. El
%     argumento va en modo matematico, sin los delimitadores ---
\newcommand{\respuesta}[1]{%
  \[
    \boxed{\;#1\;}
  \]
}

% --- resaltado inline, como pasarle el plumon a media frase. No
%     admite matematicas adentro: soul rompe con $...$, asi que una
%     formula que hay que destacar va en \respuesta ---
\sethlcolor{marcador}
\newcommand{\resaltar}[1]{\hl{#1}}

% --- justificacion al final de un renglon de \align, en chiquito y
%     entre parentesis, para decir de donde salio el paso ---
\newcommand{\razon}[1]{\quad\text{\small\itshape (#1)}}

%############################ END OF PREAMBULO.TEX #############################
%###############################################################################

%  \begin{document}
%  \portadilla
%
%  Con dos o mas profesores se declara UN solo \profesor y los nombres
%  se separan con \\, uno por renglon:
%
%  \newcommand{\profesor}{Dra. Fulana de Tal \\
%                         Dr. Mengano de Cual}
% =====================================================================


% ==================== IDIOMA Y CODIFICACIÓN ====================
\usepackage[utf8]{inputenc}
\usepackage[T1]{fontenc}
\usepackage[spanish,es-tabla,es-noshorthands]{babel}
\usepackage{lmodern}
\usepackage{microtype}

% ==================== PÁGINA ====================
\usepackage[letterpaper,margin=2.3cm,headheight=14pt]{geometry}
\usepackage{parskip}
\usepackage{fancyhdr}
\usepackage{enumitem}

% ==================== MATEMÁTICAS ====================
\usepackage{amsmath,amssymb,amsthm,mathtools}

% ==================== FIGURAS, TABLAS, CÓDIGO ====================
\usepackage{graphicx}
\usepackage{booktabs}
\usepackage{array}
% float da el especificador [H], que clava una tabla donde se
% escribio. Una traza que LaTeX manda dos paginas adelante deja
% de leerse junto al paso que la produjo
\usepackage{float}
\usepackage{xcolor}
% tikz ya viene arrastrado por tcolorbox[most], pero se carga
% explicito para no depender de ese detalle; la libreria babel
% protege a tikz de los caracteres activos del español
\usepackage{tikz}
\usetikzlibrary{babel}
% pgfplots dibuja graficas de datos (ejes, escalas log) sobre tikz
\usepackage{pgfplots}
\pgfplotsset{compat=1.18}
\usepackage{listings}
\usepackage[ruled,section]{algorithm}
\usepackage{algpseudocode}
\usepackage[most]{tcolorbox}
% soul da \hl, el unico resaltado que parte renglon; \colorbox no
\usepackage{soul}
\usepackage{hyperref}

% =====================================================================
%  ESTILO
% =====================================================================

% --- encabezado y pie de cada página ---
\pagestyle{fancy}
\fancyhf{}
\fancyhead[L]{\small\textsc{\materia}}
\fancyhead[R]{\small\periodo}
\fancyfoot[C]{\small\thepage}
\renewcommand{\headrulewidth}{0.4pt}

% --- enlaces ---
\hypersetup{
  colorlinks=true,
  linkcolor=black,
  urlcolor=blue,
  citecolor=black,
  pdftitle={Apuntes de \materia},
  pdfauthor={\alumno}
}

% --- paleta ---
\definecolor{acento}{HTML}{1F4E79}
\definecolor{fondoIdea}{HTML}{EAF2FB}
\definecolor{fondoOjo}{HTML}{FDF0E6}
\definecolor{fondoPend}{HTML}{F2F0F7}
\definecolor{comentario}{HTML}{5A7D5A}
\definecolor{cadena}{HTML}{A03030}
\definecolor{fondoCodigo}{HTML}{F7F7F7}
\definecolor{fondoResultado}{HTML}{EAF5EC}
\definecolor{marcador}{HTML}{FFF3A3}

% --- entornos de teorema; numeran por sección, o sea por clase ---
\theoremstyle{definition}
\newtheorem{definicion}{Definición}[section]
\newtheorem{ejemplo}[definicion]{Ejemplo}
\newtheorem{ejercicio}[definicion]{Ejercicio}

\theoremstyle{plain}
\newtheorem{teorema}[definicion]{Teorema}
\newtheorem{lema}[definicion]{Lema}
\newtheorem{corolario}[definicion]{Corolario}
\newtheorem{proposicion}[definicion]{Proposición}

\theoremstyle{remark}
\newtheorem*{observacion}{Observación}
\newtheorem*{quick_sum}{Resumen rápido}

% --- cajas de color para lo que se relee antes del examen ---
\newtcolorbox{idea}[1][]{
  colback=fondoIdea, colframe=acento, boxrule=0.6pt,
  arc=2pt, left=6pt, right=6pt, top=5pt, bottom=5pt,
  title={Idea clave}, fonttitle=\bfseries\small, #1
}
\newtcolorbox{ojo}[1][]{
  colback=fondoOjo, colframe=orange!70!black, boxrule=0.6pt,
  arc=2pt, left=6pt, right=6pt, top=5pt, bottom=5pt,
  title={Ojito...}, fonttitle=\bfseries\small, #1
}
\newtcolorbox{pendiente}[1][]{
  colback=fondoPend, colframe=violet!60!black, boxrule=0.6pt,
  arc=2pt, left=6pt, right=6pt, top=5pt, bottom=5pt,
  title={Pendiente por resolver}, fonttitle=\bfseries\small, #1
}
\newtcolorbox{idear}[1][]{
  colback=fondoPend, colframe=pink!60!black, boxrule=0.6pt,
  arc=2pt, left=6pt, right=6pt, top=5pt, bottom=5pt,
  title={Idea Random}, fonttitle=\bfseries\small, #1
}

% --- el resultado al que llego un ejercicio, para que se vea de
%     lejos al hojear la hoja antes del examen ---
\newtcolorbox{resultado}[1][]{
  colback=fondoResultado, colframe=green!45!black, boxrule=0.6pt,
  arc=2pt, left=6pt, right=6pt, top=5pt, bottom=5pt,
  title={Resultado}, fonttitle=\bfseries\small, #1
}

% --- código ---
\lstdefinestyle{codigo}{
  backgroundcolor=\color{fondoCodigo},
  basicstyle=\ttfamily\footnotesize,
  keywordstyle=\color{acento}\bfseries,
  commentstyle=\color{comentario}\itshape,
  stringstyle=\color{cadena},
  numbers=left, numberstyle=\tiny\color{gray}, numbersep=8pt,
  frame=single, rulecolor=\color{gray!40},
  breaklines=true, showstringspaces=false,
  tabsize=4, captionpos=b,
  extendedchars=true, inputencoding=utf8,
  literate={á}{{\'a}}1 {é}{{\'e}}1 {í}{{\'i}}1
           {ó}{{\'o}}1 {ú}{{\'u}}1 {ñ}{{\~n}}1
           {Ü}{{\"U}}1 {ü}{{\"u}}1 {¿}{{?`}}1
           {¡}{{!`}}1
}
\lstset{style=codigo}
\renewcommand{\lstlistingname}{Código}
\renewcommand{\lstlistlistingname}{Índice de código}
% --- pseudocódigo; algorithmicx arma el cuerpo y algorithm el
%     flotante que lo numera y lo pone en el índice. Las palabras
%     clave vienen en inglés de fábrica y babel no las toca, así que
%     se renombran una por una ---
\floatname{algorithm}{Algoritmo}
\renewcommand{\listalgorithmname}{Índice de algoritmos}

\algrenewcommand\algorithmicrequire{\textbf{Entrada:}}
\algrenewcommand\algorithmicensure{\textbf{Salida:}}
\algrenewcommand\algorithmicend{\textbf{fin}}
\algrenewcommand\algorithmicif{\textbf{si}}
\algrenewcommand\algorithmicthen{\textbf{entonces}}
\algrenewcommand\algorithmicelse{\textbf{si no}}
\algrenewcommand\algorithmicfor{\textbf{para}}
\algrenewcommand\algorithmicforall{\textbf{para todo}}
\algrenewcommand\algorithmicdo{\textbf{hacer}}
\algrenewcommand\algorithmicwhile{\textbf{mientras}}
\algrenewcommand\algorithmicrepeat{\textbf{repetir}}
\algrenewcommand\algorithmicuntil{\textbf{hasta que}}
\algrenewcommand\algorithmicloop{\textbf{ciclo}}
\algrenewcommand\algorithmicprocedure{\textbf{procedimiento}}
\algrenewcommand\algorithmicfunction{\textbf{función}}
\algrenewcommand\algorithmicreturn{\textbf{devolver}}
\algrenewcomment[1]{\hfill{\color{comentario}\itshape // #1}}

% Las dos que algorithmicx no trae. El \ del final es obligatorio: sin
% el TeX se come el espacio y sale "hastan-1" en vez de "hasta n-1"
\algnewcommand\Hasta{\textbf{hasta}\ }
\algnewcommand\Intercambiar{\textbf{intercambiar}\ }

% --- abreviaturas que se usan a cada rato ---
\newcommand{\N}{\mathbb{N}}
\newcommand{\Z}{\mathbb{Z}}
\newcommand{\Q}{\mathbb{Q}}
\newcommand{\R}{\mathbb{R}}
\newcommand{\C}{\mathbb{C}}
\newcommand{\abs}[1]{\left\lvert #1 \right\rvert}
\newcommand{\norm}[1]{\left\lVert #1 \right\rVert}
\newcommand{\set}[1]{\left\{ #1 \right\}}
\newcommand{\Ord}[1]{\mathcal{O}\!\left(#1\right)}
% Mayuscula la notacion grande, minuscula la chica
\newcommand{\Omg}[1]{\Omega\!\left(#1\right)}
\newcommand{\Tht}[1]{\Theta\!\left(#1\right)}
\newcommand{\ord}[1]{o\!\left(#1\right)}
\newcommand{\omg}[1]{\omega\!\left(#1\right)}
% Los tres limites sobre n, que salen en cada renglon del criterio
% del cociente para la notacion asintotica
\newcommand{\LM}{\lim_{n \to \infty}}
\newcommand{\LS}{\limsup_{n \to \infty}}
\newcommand{\LI}{\liminf_{n \to \infty}}

% --- encabezado de cada sesión: título a la izquierda, fecha a la
%     derecha. El argumento corto es el que sale en el índice ---
\newcommand{\clase}[2]{%
  \section[#1]{#1 \hfill \normalfont\small\textit{#2}}%
}

% --- portadilla y tabla de contenidos. Se llama una vez, justo
%     despues de \begin{document}. El renglon del profesor va dentro
%     del center, asi que \profesor puede traer varios nombres
%     separados por \\ y salen apilados y centrados sin tocar nada ---
\newcommand{\portadilla}{%
  \begin{center}
    {\LARGE\bfseries\materia\par}
    \vspace{0.4em}
    {\large \profesor \par}
    \vspace{0.2em}
    {\small \periodo\quad-\quad\alumno\par}
    \vspace{0.3em}
    \rule{\linewidth}{0.5pt}
  \end{center}
  \vspace{0.5em}
  \tableofcontents
  \vspace{1em}
  \rule{\linewidth}{0.5pt}%
}

% =====================================================================
%  DOCUMENTOS DE SOLUCIONES
%  Lo que usa una tanda de ejercicios resueltos. Una libreta no toca
%  nada de aqui, pero vive en el preambulo compartido para no tener
%  dos archivos que arreglar cuando cambie el estilo.
% =====================================================================

% --- encabezado compacto: titulo, subtitulo y regla. Sustituye a
%     \portadilla cuando el documento es una tanda de ejercicios: sin
%     indice, porque se hojea de corrido, no se navega ---
\newcommand{\encabezadoSoluciones}[2]{%
  % El paquete algorithm se carga con la opcion [section] para que en
  % una libreta el pseudocodigo se numere por clase. Aqui no hay
  % secciones, asi que esa numeracion sale como "Algoritmo 0.1". Se
  % corrige al vuelo, y solo en los documentos que llaman a este
  % encabezado: la libreta no se entera.
  \renewcommand{\thealgorithm}{\arabic{algorithm}}%
  \begin{center}
    {\LARGE\bfseries #1\par}
    \vspace{0.3em}
    {\small\itshape #2\par}
    \vspace{0.5em}
    \rule{\linewidth}{0.5pt}
  \end{center}
  \vspace{0.3em}%
}

% --- hoja de un solo ejercicio: un renglon chico arriba con el
%     titulo y su numero, y ya. Sin titulo grande y sin pie de pagina,
%     porque en una hoja de dos paginas cada centimetro que se va en
%     adornos es un paso del desarrollo que no cabe ---
\newcommand{\encabezadoEjercicio}[3]{%
  % Misma correccion que en \encabezadoSoluciones: sin secciones, el
  % pseudocodigo saldria numerado "0.1"
  \renewcommand{\thealgorithm}{\arabic{algorithm}}%
  \fancyhf{}%
  \fancyhead[L]{\small\textsc{#1} (ejercicio #2 de #3)}%
  % El nombre va a la derecha porque las hojas de respuestas se
  % entregan y el profesor pide que cada una lo lleve
  \fancyhead[R]{\small\alumno}%
  \renewcommand{\headrulewidth}{0.4pt}%
  \pagestyle{fancy}%
  \thispagestyle{fancy}%
}

% --- bloque tematico, numerado con romanos: I, II, III. Agrupa los
%     ejercicios que comparten tema. No dibuja regla: la del
%     \problema que sigue ya separa el titulo del primer ejercicio,
%     y dos rayas seguidas dejan un hueco feo ---
\newcounter{parte}
\renewcommand{\theparte}{\Roman{parte}}
\newcommand{\parte}[1]{%
  \bigskip
  \refstepcounter{parte}%
  \noindent{\large\bfseries \theparte. #1\par}%
}

% --- un ejercicio resuelto. Se numera solo y corrido a lo largo del
%     documento, sin reiniciar en cada parte, para poder decir "el 14"
%     sin aclarar de que bloque ---
\newcounter{problema}
\newcommand{\problema}[1]{%
  \bigskip
  \noindent\rule{\linewidth}{0.4pt}\par
  \vspace{0.3em}
  \refstepcounter{problema}%
  \noindent{\bfseries Ejercicio \theproblema\ --- #1\par}
  \vspace{0.3em}%
}

% --- rotulo del paso que sigue: Caso base, Hipotesis inductiva, Paso
%     inductivo. El texto sigue en el mismo renglon ---
\newcommand{\paso}[1]{\par\vspace{0.3em}\noindent\textbf{#1}\ }

% --- la respuesta de un ejercicio, centrada y en un cuadro. El
%     argumento va en modo matematico, sin los delimitadores ---
\newcommand{\respuesta}[1]{%
  \[
    \boxed{\;#1\;}
  \]
}

% --- resaltado inline, como pasarle el plumon a media frase. No
%     admite matematicas adentro: soul rompe con $...$, asi que una
%     formula que hay que destacar va en \respuesta ---
\sethlcolor{marcador}
\newcommand{\resaltar}[1]{\hl{#1}}

% --- justificacion al final de un renglon de \align, en chiquito y
%     entre parentesis, para decir de donde salio el paso ---
\newcommand{\razon}[1]{\quad\text{\small\itshape (#1)}}

%############################ END OF PREAMBULO.TEX #############################
%###############################################################################

%  \begin{document}
%  \portadilla
%
%  Con dos o mas profesores se declara UN solo \profesor y los nombres
%  se separan con \\, uno por renglon:
%
%  \newcommand{\profesor}{Dra. Fulana de Tal \\
%                         Dr. Mengano de Cual}
% =====================================================================


% ==================== IDIOMA Y CODIFICACIÓN ====================
\usepackage[utf8]{inputenc}
\usepackage[T1]{fontenc}
\usepackage[spanish,es-tabla,es-noshorthands]{babel}
\usepackage{lmodern}
\usepackage{microtype}

% ==================== PÁGINA ====================
\usepackage[letterpaper,margin=2.3cm,headheight=14pt]{geometry}
\usepackage{parskip}
\usepackage{fancyhdr}
\usepackage{enumitem}

% ==================== MATEMÁTICAS ====================
\usepackage{amsmath,amssymb,amsthm,mathtools}

% ==================== FIGURAS, TABLAS, CÓDIGO ====================
\usepackage{graphicx}
\usepackage{booktabs}
\usepackage{array}
% float da el especificador [H], que clava una tabla donde se
% escribio. Una traza que LaTeX manda dos paginas adelante deja
% de leerse junto al paso que la produjo
\usepackage{float}
\usepackage{xcolor}
% tikz ya viene arrastrado por tcolorbox[most], pero se carga
% explicito para no depender de ese detalle; la libreria babel
% protege a tikz de los caracteres activos del español
\usepackage{tikz}
\usetikzlibrary{babel}
% pgfplots dibuja graficas de datos (ejes, escalas log) sobre tikz
\usepackage{pgfplots}
\pgfplotsset{compat=1.18}
\usepackage{listings}
\usepackage[ruled,section]{algorithm}
\usepackage{algpseudocode}
\usepackage[most]{tcolorbox}
% soul da \hl, el unico resaltado que parte renglon; \colorbox no
\usepackage{soul}
\usepackage{hyperref}

% =====================================================================
%  ESTILO
% =====================================================================

% --- encabezado y pie de cada página ---
\pagestyle{fancy}
\fancyhf{}
\fancyhead[L]{\small\textsc{\materia}}
\fancyhead[R]{\small\periodo}
\fancyfoot[C]{\small\thepage}
\renewcommand{\headrulewidth}{0.4pt}

% --- enlaces ---
\hypersetup{
  colorlinks=true,
  linkcolor=black,
  urlcolor=blue,
  citecolor=black,
  pdftitle={Apuntes de \materia},
  pdfauthor={\alumno}
}

% --- paleta ---
\definecolor{acento}{HTML}{1F4E79}
\definecolor{fondoIdea}{HTML}{EAF2FB}
\definecolor{fondoOjo}{HTML}{FDF0E6}
\definecolor{fondoPend}{HTML}{F2F0F7}
\definecolor{comentario}{HTML}{5A7D5A}
\definecolor{cadena}{HTML}{A03030}
\definecolor{fondoCodigo}{HTML}{F7F7F7}
\definecolor{fondoResultado}{HTML}{EAF5EC}
\definecolor{marcador}{HTML}{FFF3A3}

% --- entornos de teorema; numeran por sección, o sea por clase ---
\theoremstyle{definition}
\newtheorem{definicion}{Definición}[section]
\newtheorem{ejemplo}[definicion]{Ejemplo}
\newtheorem{ejercicio}[definicion]{Ejercicio}

\theoremstyle{plain}
\newtheorem{teorema}[definicion]{Teorema}
\newtheorem{lema}[definicion]{Lema}
\newtheorem{corolario}[definicion]{Corolario}
\newtheorem{proposicion}[definicion]{Proposición}

\theoremstyle{remark}
\newtheorem*{observacion}{Observación}
\newtheorem*{quick_sum}{Resumen rápido}

% --- cajas de color para lo que se relee antes del examen ---
\newtcolorbox{idea}[1][]{
  colback=fondoIdea, colframe=acento, boxrule=0.6pt,
  arc=2pt, left=6pt, right=6pt, top=5pt, bottom=5pt,
  title={Idea clave}, fonttitle=\bfseries\small, #1
}
\newtcolorbox{ojo}[1][]{
  colback=fondoOjo, colframe=orange!70!black, boxrule=0.6pt,
  arc=2pt, left=6pt, right=6pt, top=5pt, bottom=5pt,
  title={Ojito...}, fonttitle=\bfseries\small, #1
}
\newtcolorbox{pendiente}[1][]{
  colback=fondoPend, colframe=violet!60!black, boxrule=0.6pt,
  arc=2pt, left=6pt, right=6pt, top=5pt, bottom=5pt,
  title={Pendiente por resolver}, fonttitle=\bfseries\small, #1
}
\newtcolorbox{idear}[1][]{
  colback=fondoPend, colframe=pink!60!black, boxrule=0.6pt,
  arc=2pt, left=6pt, right=6pt, top=5pt, bottom=5pt,
  title={Idea Random}, fonttitle=\bfseries\small, #1
}

% --- el resultado al que llego un ejercicio, para que se vea de
%     lejos al hojear la hoja antes del examen ---
\newtcolorbox{resultado}[1][]{
  colback=fondoResultado, colframe=green!45!black, boxrule=0.6pt,
  arc=2pt, left=6pt, right=6pt, top=5pt, bottom=5pt,
  title={Resultado}, fonttitle=\bfseries\small, #1
}

% --- código ---
\lstdefinestyle{codigo}{
  backgroundcolor=\color{fondoCodigo},
  basicstyle=\ttfamily\footnotesize,
  keywordstyle=\color{acento}\bfseries,
  commentstyle=\color{comentario}\itshape,
  stringstyle=\color{cadena},
  numbers=left, numberstyle=\tiny\color{gray}, numbersep=8pt,
  frame=single, rulecolor=\color{gray!40},
  breaklines=true, showstringspaces=false,
  tabsize=4, captionpos=b,
  extendedchars=true, inputencoding=utf8,
  literate={á}{{\'a}}1 {é}{{\'e}}1 {í}{{\'i}}1
           {ó}{{\'o}}1 {ú}{{\'u}}1 {ñ}{{\~n}}1
           {Ü}{{\"U}}1 {ü}{{\"u}}1 {¿}{{?`}}1
           {¡}{{!`}}1
}
\lstset{style=codigo}
\renewcommand{\lstlistingname}{Código}
\renewcommand{\lstlistlistingname}{Índice de código}
% --- pseudocódigo; algorithmicx arma el cuerpo y algorithm el
%     flotante que lo numera y lo pone en el índice. Las palabras
%     clave vienen en inglés de fábrica y babel no las toca, así que
%     se renombran una por una ---
\floatname{algorithm}{Algoritmo}
\renewcommand{\listalgorithmname}{Índice de algoritmos}

\algrenewcommand\algorithmicrequire{\textbf{Entrada:}}
\algrenewcommand\algorithmicensure{\textbf{Salida:}}
\algrenewcommand\algorithmicend{\textbf{fin}}
\algrenewcommand\algorithmicif{\textbf{si}}
\algrenewcommand\algorithmicthen{\textbf{entonces}}
\algrenewcommand\algorithmicelse{\textbf{si no}}
\algrenewcommand\algorithmicfor{\textbf{para}}
\algrenewcommand\algorithmicforall{\textbf{para todo}}
\algrenewcommand\algorithmicdo{\textbf{hacer}}
\algrenewcommand\algorithmicwhile{\textbf{mientras}}
\algrenewcommand\algorithmicrepeat{\textbf{repetir}}
\algrenewcommand\algorithmicuntil{\textbf{hasta que}}
\algrenewcommand\algorithmicloop{\textbf{ciclo}}
\algrenewcommand\algorithmicprocedure{\textbf{procedimiento}}
\algrenewcommand\algorithmicfunction{\textbf{función}}
\algrenewcommand\algorithmicreturn{\textbf{devolver}}
\algrenewcomment[1]{\hfill{\color{comentario}\itshape // #1}}

% Las dos que algorithmicx no trae. El \ del final es obligatorio: sin
% el TeX se come el espacio y sale "hastan-1" en vez de "hasta n-1"
\algnewcommand\Hasta{\textbf{hasta}\ }
\algnewcommand\Intercambiar{\textbf{intercambiar}\ }

% --- abreviaturas que se usan a cada rato ---
\newcommand{\N}{\mathbb{N}}
\newcommand{\Z}{\mathbb{Z}}
\newcommand{\Q}{\mathbb{Q}}
\newcommand{\R}{\mathbb{R}}
\newcommand{\C}{\mathbb{C}}
\newcommand{\abs}[1]{\left\lvert #1 \right\rvert}
\newcommand{\norm}[1]{\left\lVert #1 \right\rVert}
\newcommand{\set}[1]{\left\{ #1 \right\}}
\newcommand{\Ord}[1]{\mathcal{O}\!\left(#1\right)}
% Mayuscula la notacion grande, minuscula la chica
\newcommand{\Omg}[1]{\Omega\!\left(#1\right)}
\newcommand{\Tht}[1]{\Theta\!\left(#1\right)}
\newcommand{\ord}[1]{o\!\left(#1\right)}
\newcommand{\omg}[1]{\omega\!\left(#1\right)}
% Los tres limites sobre n, que salen en cada renglon del criterio
% del cociente para la notacion asintotica
\newcommand{\LM}{\lim_{n \to \infty}}
\newcommand{\LS}{\limsup_{n \to \infty}}
\newcommand{\LI}{\liminf_{n \to \infty}}

% --- encabezado de cada sesión: título a la izquierda, fecha a la
%     derecha. El argumento corto es el que sale en el índice ---
\newcommand{\clase}[2]{%
  \section[#1]{#1 \hfill \normalfont\small\textit{#2}}%
}

% --- portadilla y tabla de contenidos. Se llama una vez, justo
%     despues de \begin{document}. El renglon del profesor va dentro
%     del center, asi que \profesor puede traer varios nombres
%     separados por \\ y salen apilados y centrados sin tocar nada ---
\newcommand{\portadilla}{%
  \begin{center}
    {\LARGE\bfseries\materia\par}
    \vspace{0.4em}
    {\large \profesor \par}
    \vspace{0.2em}
    {\small \periodo\quad-\quad\alumno\par}
    \vspace{0.3em}
    \rule{\linewidth}{0.5pt}
  \end{center}
  \vspace{0.5em}
  \tableofcontents
  \vspace{1em}
  \rule{\linewidth}{0.5pt}%
}

% =====================================================================
%  DOCUMENTOS DE SOLUCIONES
%  Lo que usa una tanda de ejercicios resueltos. Una libreta no toca
%  nada de aqui, pero vive en el preambulo compartido para no tener
%  dos archivos que arreglar cuando cambie el estilo.
% =====================================================================

% --- encabezado compacto: titulo, subtitulo y regla. Sustituye a
%     \portadilla cuando el documento es una tanda de ejercicios: sin
%     indice, porque se hojea de corrido, no se navega ---
\newcommand{\encabezadoSoluciones}[2]{%
  % El paquete algorithm se carga con la opcion [section] para que en
  % una libreta el pseudocodigo se numere por clase. Aqui no hay
  % secciones, asi que esa numeracion sale como "Algoritmo 0.1". Se
  % corrige al vuelo, y solo en los documentos que llaman a este
  % encabezado: la libreta no se entera.
  \renewcommand{\thealgorithm}{\arabic{algorithm}}%
  \begin{center}
    {\LARGE\bfseries #1\par}
    \vspace{0.3em}
    {\small\itshape #2\par}
    \vspace{0.5em}
    \rule{\linewidth}{0.5pt}
  \end{center}
  \vspace{0.3em}%
}

% --- hoja de un solo ejercicio: un renglon chico arriba con el
%     titulo y su numero, y ya. Sin titulo grande y sin pie de pagina,
%     porque en una hoja de dos paginas cada centimetro que se va en
%     adornos es un paso del desarrollo que no cabe ---
\newcommand{\encabezadoEjercicio}[3]{%
  % Misma correccion que en \encabezadoSoluciones: sin secciones, el
  % pseudocodigo saldria numerado "0.1"
  \renewcommand{\thealgorithm}{\arabic{algorithm}}%
  \fancyhf{}%
  \fancyhead[L]{\small\textsc{#1} (ejercicio #2 de #3)}%
  % El nombre va a la derecha porque las hojas de respuestas se
  % entregan y el profesor pide que cada una lo lleve
  \fancyhead[R]{\small\alumno}%
  \renewcommand{\headrulewidth}{0.4pt}%
  \pagestyle{fancy}%
  \thispagestyle{fancy}%
}

% --- bloque tematico, numerado con romanos: I, II, III. Agrupa los
%     ejercicios que comparten tema. No dibuja regla: la del
%     \problema que sigue ya separa el titulo del primer ejercicio,
%     y dos rayas seguidas dejan un hueco feo ---
\newcounter{parte}
\renewcommand{\theparte}{\Roman{parte}}
\newcommand{\parte}[1]{%
  \bigskip
  \refstepcounter{parte}%
  \noindent{\large\bfseries \theparte. #1\par}%
}

% --- un ejercicio resuelto. Se numera solo y corrido a lo largo del
%     documento, sin reiniciar en cada parte, para poder decir "el 14"
%     sin aclarar de que bloque ---
\newcounter{problema}
\newcommand{\problema}[1]{%
  \bigskip
  \noindent\rule{\linewidth}{0.4pt}\par
  \vspace{0.3em}
  \refstepcounter{problema}%
  \noindent{\bfseries Ejercicio \theproblema\ --- #1\par}
  \vspace{0.3em}%
}

% --- rotulo del paso que sigue: Caso base, Hipotesis inductiva, Paso
%     inductivo. El texto sigue en el mismo renglon ---
\newcommand{\paso}[1]{\par\vspace{0.3em}\noindent\textbf{#1}\ }

% --- la respuesta de un ejercicio, centrada y en un cuadro. El
%     argumento va en modo matematico, sin los delimitadores ---
\newcommand{\respuesta}[1]{%
  \[
    \boxed{\;#1\;}
  \]
}

% --- resaltado inline, como pasarle el plumon a media frase. No
%     admite matematicas adentro: soul rompe con $...$, asi que una
%     formula que hay que destacar va en \respuesta ---
\sethlcolor{marcador}
\newcommand{\resaltar}[1]{\hl{#1}}

% --- justificacion al final de un renglon de \align, en chiquito y
%     entre parentesis, para decir de donde salio el paso ---
\newcommand{\razon}[1]{\quad\text{\small\itshape (#1)}}

%############################ END OF PREAMBULO.TEX #############################
%###############################################################################


% =====================================================================
\begin{document}
% =====================================================================

\portadilla
\newpage

% =====================================================================
%  Clase 1 - Intro al curso 
% =====================================================================

\clase{Introducción}{17 de agosto de 2026}

\begin{quick_sum}
    En esta sesión vimos una intro... que es lo que se espera de nosotros en la
    materia, cómo se va a evaluar y qué temas se van a ver. Nos dijeron que
    no es tanto como para hacer demostraciones sino como para entender los 
    conceptos y saber usar esas herramientas en la computación.
\end{quick_sum}

\newpage

% =====================================================================
%  Clase 2 - Repaso 
% =====================================================================

\clase{Repaso}{19 de agosto de 2026}

\begin{quick_sum}
    En esta clase iniciamos con el repaso de algunos conceptos
    matemáticos en los que se basa la materia y en general las mates
    para compu. Los conceptos que vimos son el de algebra, conjuntos,
    subconjuntos, operaciones con conjuntos y producto cartesiano.
\end{quick_sum}


\subsection*{Álgebra}

El álgebra es la parte de las matemáticas que estudia las operaciones,
sus signos y las leyes que las rigen.

% =====================================================================

\subsection*{Conjuntos}

Los conjuntos permiten agrupar elementos que comparten algo y operar
entre esos grupos. Son la forma más básica de describir qué objetos
estamos estudiando.

\begin{definicion}[Conjunto]
    Colección bien definida de objetos que comparten alguna
    característica. Puede ser finito o infinito.
\end{definicion}

Los elementos de un conjunto son distintos entre sí. Si al escribirlo
se repite uno, cuenta una sola vez. Los enteros son el ejemplo de
conjunto infinito.

\subsubsection*{Notación}

Un conjunto se describe por comprensión dando la regla que cumplen sus
elementos:
\[
    X = \set{x \mid \text{regla que cumple } x}
\]

O por extensión listando sus elementos:
\[
    X = \set{x_1, x_2, x_3, \ldots, x_n}
\]

Los tres símbolos que se usan a cada rato:

\begin{itemize}[itemsep=0.1em]
    \item \textbf{Pertenencia} ($x \in X$): dice que $x$ es uno de los
          elementos de $X$.
    \item \textbf{Cardinalidad} ($\abs{X}$): el número de elementos que
          tiene el conjunto.
    \item \textbf{Universo} ($U$): todas las posibilidades sobre las
          que se define la regla. Sin fijar $U$, la regla no delimita
          nada.
\end{itemize}

\subsubsection*{Conjunto vacío e igualdad}

El conjunto vacío $\emptyset$ no tiene elementos, su cardinalidad es
$0$ y es subconjunto de cualquier conjunto.

Dos conjuntos son iguales cuando lo son elemento a elemento.

\subsubsection*{Subconjuntos}

$A \subseteq B$ cuando todo elemento de $A$ está también en $B$. Todo
conjunto es subconjunto de sí mismo, $A \subseteq A$.

El subconjunto es \textbf{propio}, $A \subset B$, cuando además
$A \neq B$: la contención se cumple pero $B$ tiene al menos un elemento
que $A$ no tiene.

\subsubsection*{Ejercicio clase}
La Dra dejo un ejercicio de que así:

Sea
\[
    \mathcal{U} = \set{1, 2, 3, 4, \set{1, 2}, \set{3, 5},
                       \set{1, 2, 3, 4}}
    \qquad
    A = \set{1, 2, 3, 4}
\]
Determinar si lo siguiente es verdadero o falso:
\begin{enumerate}[label=\alph*)]
    \item $\abs{\mathcal{U}} = 7$ \textbf{Verdadero}, pq tiene 7
          elementos distintos
    \item $A \subseteq \mathcal{U}$ \textbf{Verdadero}, pq todos los
          elementos de $A$ están en $\mathcal{U}$
    \item $A \in \mathcal{U}$ \textbf{Verdadero}, pq $A$ es un elemento
          de $\mathcal{U}$
    \item $\abs{A} = 4$ \textbf{Verdadero}, pq tiene 4 elementos
          distintos
\end{enumerate}

\subsubsection*{Conjunto potencia}

\begin{definicion}[Conjunto potencia]
    El conjunto formado por todos los subconjuntos de $A$, incluidos
    $\emptyset$ y el propio $A$. Se escribe $\mathcal{P}(A)$.
\end{definicion}

Su cardinalidad es $\abs{\mathcal{P}(A)} = 2^{\abs{A}}$. De dónde sale
el $2^n$: armar un subconjunto es decidir, para cada uno de los $n$
elementos, si entra o no entra. Son dos opciones independientes por
elemento, así que salen $2 \times 2 \times \cdots \times 2 = 2^n$
combinaciones, y cada una da un subconjunto distinto.

\begin{ejemplo}
    Con $A = \set{1, 2}$ hay $n = 2$ elementos y $2^2 = 4$
    subconjuntos:
    \[
        \mathcal{P}(A) = \set{\ \emptyset,\ \set{1},\ \set{2},\
                              \set{1, 2}\ }
    \]
\end{ejemplo}

% =====================================================================

\subsection*{Operaciones con conjuntos}

Todas las que vimos son \textbf{binarias}: toman dos operandos.

\subsubsection*{Cerradura}

Si $a$ y $b$ pertenecen a los enteros y el resultado de operarlos
también es un entero, la operación es \textbf{cerrada} en los enteros.
La cerradura sirve para ponerle etiqueta y nombre al álgebra en la que
se está trabajando.

\subsubsection*{Las cinco operaciones}

\begin{table}[htbp]
    \centering
    \begin{tabular}{llp{6cm}}
        \toprule
        \textbf{Operación} & \textbf{Símbolo} & \textbf{Qué devuelve}\\
        \midrule
        Unión                & $A \cup B$
          & Lo que está en $A$, en $B$ o en ambos \\
        Intersección         & $A \cap B$
          & Lo que está en los dos a la vez \\
        Diferencia           & $A \setminus B$
          & Lo de $A$ que no está en $B$ \\
        Diferencia simétrica & $A \mathbin{\triangle} B$
          & Lo que está en uno de los dos, pero no en ambos \\
        Complemento          & $A^{c}$
          & Lo de $U$ que no está en $A$ \\
        \bottomrule
    \end{tabular}
    \caption{Las cinco operaciones vistas en clase.}
    \label{tab:operaciones}
\end{table}

% =====================================================================

\subsection*{Producto cartesiano}

\begin{definicion}[Producto cartesiano]
    $A \times B$ es el conjunto de todos los pares ordenados cuyo
    primer elemento sale de $A$ y cuyo segundo sale de $B$:
    \[
        A \times B = \set{(a, b) \mid a \in A,\ b \in B}
    \]
\end{definicion}

Se emparejan todos los elementos de $A$ con todos los de $B$, así que
la cardinalidad es $\abs{A \times B} = \abs{A} \cdot \abs{B}$.

\newpage

% =====================================================================
%  Clase 3 - Relaciones y funciones y Álgebras
% =====================================================================

\clase{Relaciones y funciones}{24 de agosto de 2026}

\begin{quick_sum}
    En esta clase vimos un repaso de relaciones y funciones, tambien vimos
    una intro a las álgebras.
\end{quick_sum}

\subsection*{Relaciones}

\begin{definicion}[Relación binaria]
    Una relación binaria $R$ de un conjunto $X$ en un conjunto $Y$ es
    un subconjunto del producto cartesiano $X \times Y$.
\end{definicion}

Si $(x, y) \in R$, escribimos $xRy$ para denotar que $x$ está
relacionada con $y$. Si $X = Y$ decimos que $R$ es una relación
binaria \textbf{sobre} $X$.

Como una relación es un subconjunto de $X \times Y$, todo lo de
conjuntos aplica: $\emptyset$ es una relación, y hay
$2^{\abs{X} \cdot \abs{Y}}$ relaciones distintas entre $X$ y $Y$.

\subsubsection*{Dominio, codominio y rango}

El \textbf{dominio} está formado por los elementos del primer conjunto
del que estamos hablando, y el \textbf{codominio} por los del segundo.
El \textbf{rango} o imagen son los elementos que sí se usan, o sea los
que participan de forma efectiva:
\[
    \mathrm{Dom}(R) = \set{x \in X \mid \exists\, y \in Y,\ (x,y) \in R}
    \qquad
    \mathrm{Ran}(R) = \set{y \in Y \mid \exists\, x \in X,\ (x,y) \in R}
\]

La diferencia entre codominio y rango es que el codominio se declara y
el rango se calcula: $Y$ es lo que la relación podría alcanzar,
$\mathrm{Ran}(R)$ es lo que de verdad alcanza.

\begin{ejemplo}
    Sean $X = \set{1, 2, 3}$ y $Y = \set{a, b}$, con
    $R = \set{(1,a),\ (1,b),\ (3,a)}$. Entonces
    $\mathrm{Dom}(R) = \set{1, 3}$ y $\mathrm{Ran}(R) = \set{a, b}$.
    El $2$ no participa y el codominio sigue siendo todo $Y$.
\end{ejemplo}

\subsubsection*{Relación n-aria}

Una relación n-aria es un subconjunto del producto cartesiano de $n$
conjuntos $A_1 \times A_2 \times \cdots \times A_n$. Esos son los
conjuntos dominio y $n$ es el \textbf{grado} de la relación.

$R$ consiste en un conjunto de n-tuplas $(a_1, a_2, a_3, \ldots, a_n)$
tales que el $i$-ésimo elemento de la tupla proviene del conjunto
$A_i$. La relación binaria es el caso $n = 2$.

\subsubsection*{Propiedades de las relaciones}

Se definen sobre una relación $R$ en un conjunto $X$.

\begin{itemize}[itemsep=0.2em]
    \item \textbf{Reflexividad}: para todo $x$, la tupla $(x,x)$ está
          en la relación. Es decir, $x$ se relaciona consigo mismo
          para todo $x \in X$.
    \item \textbf{Simetría}: para todo $x, y \in X$, si $(x,y)$ está
          en la relación, entonces $(y,x)$ también.
    \item \textbf{Antisimetría}: para todo $x, y \in X$, si $(x,y)$ y
          $(y,x)$ están en la relación, entonces $x = y$. Dicho al
          revés, que es más fácil de ver: los únicos pares que pueden
          ir y venir son los de un elemento consigo mismo.
    \item \textbf{Transitividad}: para todo $x, y, z \in X$, si $xRy$
          y $yRz$, entonces $xRz$.
\end{itemize}

\begin{ojo}
    Simetría y antisimetría no son opuestas, se pueden cumplir las dos
    a la vez. La igualdad es el caso: si $x = y$ entonces $y = x$
    (simétrica), y si ambas se cumplen entonces $x = y$
    (antisimétrica).
\end{ojo}

% =====================================================================

\subsection*{Funciones}

Son un tipo especial de relaciones. Es una relación pero con la
restricción de que cada elemento de $A$ debe tener un elemento de
salida, o sea que solo debe aparecer una vez en el dominio.

\begin{definicion}[Función]
    $f: A \to B$ es una relación de $A$ en $B$ tal que para cada
    $a \in A$ existe exactamente un $b \in B$ con $(a,b) \in f$.
\end{definicion}

\textbf{Mapear} es una función: un mapa es una forma de asociar
objetos únicos a cada elemento de un conjunto dado.

\subsubsection*{Tipos de funciones}

\begin{itemize}[itemsep=0.2em]
    \item \textbf{Inyectiva}: la relación se da $1$ a $1$, un elemento
          de $A$ no se puede relacionar con dos de $B$. Formalmente,
          si $f(x) = f(y)$ entonces $x = y$.
    \item \textbf{Sobreyectiva}: todos los elementos de $B$ participan,
          ninguno se queda solo. O sea, el rango es todo el codominio.
    \item \textbf{Biyectiva}: ambas, inyectiva y sobreyectiva. Rango
          igual a codominio, y cada elemento de $B$ le toca a exactamente
          un elemento de $A$.
\end{itemize}

La biyección es la que importa porque es la única que se puede
invertir: existe $f^{-1}: B \to A$ y también es función.

Dado todo lo que hemos visto hasta ahora, todo eso se define dentro de
una estructura algebraica.

\newpage

\clase{Algebras}{24 de agosto de 2026}

\subsection*{Estructura algebraica}

\begin{definicion}[Estructura sobre un conjunto]
    Conjunto de funciones, reglas y restricciones que dan significado a
    una colección de objetos.
\end{definicion}

Se dice que hay una \textbf{composición} en un conjunto no vacío $S$
si a cada par de elementos de $S$ se le asigna, de cierta forma y de
manera única, un elemento de $S$.

\subsubsection*{Aridad}

De acuerdo al número de argumentos requeridos para que un operador
pueda realizar el cálculo (aridad), las operaciones pueden ser:

\begin{itemize}[itemsep=0.05em]
    \item Nularias
    \item Unarias
    \item Binarias
    \item Ternarias
    \item N-arias
\end{itemize}

Las operaciones algebraicas básicas son la \textbf{adición} y la
\textbf{multiplicación}.

\subsubsection*{Propiedades de una operación binaria}

Sea $*$ una operación binaria sobre un conjunto $S$.

\begin{itemize}[itemsep=0.2em]
    \item \textbf{Cerradura}: $a * b$ sigue estando en $S$. En los
          naturales la suma es cerrada, pero la resta no lo es:
          $3 - 5$ ya no es natural.
    \item \textbf{Asociatividad}: $(a * b) * c = a * (b * c)$, o sea
          que el orden de agrupar no cambia el resultado.
    \item \textbf{Identidad}: existe $e \in S$ tal que
          $a * e = e * a = a$. El $0$ para la suma, el $1$ para el
          producto.
    \item \textbf{Conmutatividad}: $a * b = b * a$.
    \item \textbf{Distributividad}: relaciona dos operaciones,
          $a * (b + c) = (a * b) + (a * c)$.
    \item \textbf{Idempotencia}: $a * a = a$. La unión y la
          intersección de conjuntos la cumplen.
\end{itemize}

\begin{idea}
    La cerradura es la que decide de qué álgebra estamos hablando. Es
    la misma idea que ya salió con los conjuntos: si la operación se
    sale del conjunto, el conjunto no alcanza para describir lo que se
    está haciendo.
\end{idea}

\newpage

% =====================================================================
%  Clase 26 Ago - Continuacion de estructuras algebraicas
% =====================================================================

\clase{Continuación (Estructuras algebraicas)}{26 de agosto de 2026}

\subsection*{Definición y para qué sirven}

\begin{definicion}[Estructura algebraica]
    Un conjunto no vacío $S$ es llamado estructura algebraica si al
    menos una composición, o más exactamente una $n$-tupla de
    composiciones, está definida dentro de ella.
\end{definicion}

Sirven para diferenciar qué elementos se están utilizando y qué
operaciones son válidas entre ellos. Por ejemplo, si el conjunto $S$
contiene a los enteros, no se puede aplicar un producto vectorial.

Una estructura algebraica está compuesta por un conjunto $S$ y las
operaciones que pueden aplicarse sobre $S$, que satisfacen ciertas
propiedades:

\begin{itemize}[itemsep=0.05em]
    \item Operación cerrada o interna (clausura)
    \item Asociativa
    \item Conmutativa
    \item Elemento neutro
    \item Elemento inverso (simetría)
    \item Distributividad
\end{itemize}

El elemento inverso es el que no había salido antes: para cada
$a \in S$ existe $a^{-1} \in S$ tal que $a * a^{-1} = e$, con $e$ el
neutro. En $(\Z, +)$ el inverso de $3$ es $-3$; en $(\Z, \cdot)$ el
$3$ no tiene inverso, porque $1/3$ ya no es entero.

% =====================================================================

\subsection*{Estructuras con una operación}

\subsubsection*{Magma}

Hace referencia a un conjunto equipado con un operador binario. Es la
estructura algebraica más básica: consta de un conjunto $M$ y una
única operación binaria cerrada, sin que se imponga ninguna otra
propiedad.

\subsubsection*{Semigrupo}

Es un conjunto no vacío en el cual está definida una operación
binaria asociativa, $(S, *)$.

\subsubsection*{Monoide}

Estructura algebraica formada por un conjunto no vacío $S$ y una
operación binaria asociativa que además tiene elemento neutro.

\begin{ejemplo}
    $(\N, +)$ es asociativa y su elemento neutro es el $0$, así que
    es un monoide.

    $(\N, \cdot)$ es asociativa y su neutro es el $1$, así que
    también es un monoide.

    $(\Z^{+}, +)$ es asociativa pero no tiene neutro, porque el $0$
    no está en $\Z^{+}$, así que solo llega a semigrupo.
\end{ejemplo}

\subsubsection*{Grupo}

Si $S$ es un conjunto no vacío y $*$ una operación binaria en $S$,
entonces $(S, *)$ es un grupo si cumple con ser asociativo, tener
elemento neutro y tener inversos.

\begin{idea}
    Dicho de otro modo, un grupo es un monoide en el que cada
    elemento tiene un inverso. Cada estructura de la lista es la
    anterior más una propiedad, y por eso el orden en que se
    presentan no es arbitrario.
\end{idea}

\subsubsection*{Grupo abeliano o conmutativo}

Es una estructura de tipo grupo que además cumple con ser
conmutativa.

% =====================================================================

\subsection*{Estructuras con dos operaciones}

\subsubsection*{Anillo}

Es una estructura algebraica formada por un conjunto no vacío $S$ y
dos operaciones internas, llamadas usualmente suma y producto,
$(S, +, \cdot)$, que cumplen ciertas propiedades.

Respecto a la operación suma es cerrado, asociativo, tiene elemento
neutro, tiene elemento inverso y es conmutativo. Respecto a la
operación producto es cerrado, asociativo, distributivo y
conmutativo.

\begin{ojo}
    Lo del producto conmutativo hay que revisarlo. En la tabla de
    abajo, que es la que dio el profesor, la conmutatividad del
    producto es justo lo que separa al anillo del anillo conmutativo:
    el anillo la tiene solo en la suma. Es probable que esa última
    propiedad de la lista corresponda al anillo conmutativo.
\end{ojo}

\subsubsection*{Tabla de todas las estructuras}

En la tabla, \textbf{S} quiere decir que la propiedad se cumple
respecto a la suma y \textbf{M} respecto a la multiplicación. Las
primeras cuatro son magmas, o sea de una sola operación, y las
últimas cuatro son semianillos.

Cada renglón de la tabla es el anterior más una propiedad:

\begin{itemize}[itemsep=0.2em]
    \item \textbf{Semianillo}: las dos operaciones son asociativas y
          el producto distribuye sobre la suma.
    \item \textbf{Subanillo}: además la suma tiene neutro e inverso.
    \item \textbf{Anillo}: además la suma es conmutativa.
    \item \textbf{Anillo conmutativo}: además el producto también es
          conmutativo.
    \item \textbf{Campo}: además el producto tiene neutro e inverso.
\end{itemize}

\begin{table}[htbp]
  \centering
  \footnotesize
  \begin{tabular}{lccccc}
    \toprule
    \textbf{Magma} & \textbf{Asociativo} & \textbf{Neutro ($e$)}
      & \textbf{Inverso} & \textbf{Conmutativo}
      & \textbf{Distributivo} \\
    \midrule
    Semigrupo           & X   &     &     &     &   \\
    Monoide             & X   & X   &     &     &   \\
    Grupo               & X   & X   & X   &     &   \\
    Grupo abeliano      & X   & X   & X   & X   &   \\
    \midrule
    \textbf{Semianillo} & S,M &     &     &     & M \\
    \midrule
    Subanillo           & S,M & S   & S   &     & M \\
    Anillo              & S,M & S   & S   & S   & M \\
    Anillo conmutativo  & S,M & S   & S   & S,M & M \\
    Campo               & S,M & S,M & S,M & S,M & M \\
    \bottomrule
  \end{tabular}
  \caption{Propiedades que cumple cada estructura. Todas son
    operaciones cerradas.}
  \label{tab:estructuras}
\end{table}



\end{document}
%############################## END OF LIBRETA.TEX #############################
%###############################################################################
